\subsection{Probabilistic Membership Testing in Distributed Systems}

We consider the dictionary problem, which consists of maintaining a dynamic set $S \subseteq \mathcal{U}$ and supporting membership queries. This capability is a fundamental primitive for distributed protocols concerned with set reconciliation, the process of computing the symmetric difference between sets held on disparate nodes. In a large-scale distributed environment, classical deterministic solutions to the membership problem, such as hash tables, impose space requirements proportional to both the number of elements and their individual sizes. For $n$ elements of size $|x|$ bits, a hash table requires $\Theta(n \cdot |x|)$ space. When data is partitioned across nodes, the communication cost of synchronizing these structures scales with the same factor. For large objects or high-cardinality sets, this cost renders deterministic approaches architecturally infeasible in bandwidth-constrained systems \cite{Liu2014ParBFAP, Chang2006BigtableAD}.

Probabilistic data structures resolve this constraint by trading a quantifiable measure of precision for substantial gains in space and communication efficiency. A Bloom filter, the canonical instance of this class, requires $\Theta(n \log(1/\varepsilon))$ space independent of element size, where $\varepsilon$ denotes the target false positive rate. This sub-linear space complexity enables the design of systems that prioritize performance and availability over strict consistency. We establish the formal properties of probabilistic membership oracles and analyze the specific constraints that distributed deployment imposes on their design.

\subsubsection{Membership Queries in Large-Scale Systems}

A one-sided error membership oracle guarantees correctness for negative answers, a property of paramount importance in distributed systems. A definitive negative answer allows a system to avoid expensive, high-latency operations, such as disk I/O or network round-trips, which are often the primary bottlenecks to performance \cite{Tarkoma2012TheoryAP, Chang2006BigtableAD}. For instance, a storage system may consult such an oracle to determine if a specific data block might exist on a remote disk before initiating a costly read operation; a negative response obviates the need for any disk access \cite{Chang2006BigtableAD, Tarkoma2012TheoryAP}.

\begin{definition}[One-Sided Error Membership Oracle]
    A data structure $\mathcal{D}$ representing a set $S$ is a one-sided error membership oracle with a false positive rate $\varepsilon$ if for any element $x \in S$, the probability that $\mathcal{D}.\textsc{Query}(x)$ returns true is 1, and for any element $y \notin S$, the probability that $\mathcal{D}.\textsc{Query}(y)$ returns true is at most $\varepsilon$.
\end{definition}

By definition, the oracle never produces false negatives. The adoption of such a structure is not an ad-hoc optimization but a principled architectural strategy. It represents a deliberate trade-off where a system prioritizes lower latency over strict consistency. By accepting a bounded and well-defined error probability, these algorithms enable a level of scalability and performance that would be unattainable with exact, deterministic methods \cite{Tarkoma2012TheoryAP, Chang2006BigtableAD}.

\subsubsection{The Standard Bloom Filter}

The Bloom filter is the canonical implementation of a one-sided error membership oracle \cite{Tarkoma2012TheoryAP}. It achieves its space efficiency by storing compact representations of elements, derived from hash functions, rather than the elements themselves.

\begin{definition}[Bloom Filter \cite{Tarkoma2012TheoryAP, Broder2004NetworkAO}]
    A Bloom filter for a set $S$ with expected cardinality $|S| \leq n$ from a universe $\mathcal{U}$ consists of a bit array $B$ of length $m$ initialized to all zeros, and a collection of $k$ independent hash functions $\mathcal{H} = \{h_1, \ldots, h_k\}$, where each $h_i : \mathcal{U} \to \{0, 1, \ldots, m-1\}$.
\end{definition}

To insert an element $x \in S$, we set the bits at positions $B[h_i(x)]$ to 1 for all $i \in \{1, \ldots, k\}$. To query for an element $y$, the operation returns true if and only if $B[h_i(y)] = 1$ for all $i \in \{1, \ldots, k\}$. A query for any element $y \in S$ is guaranteed to return true, as all corresponding bits will have been set during its insertion. This structure therefore satisfies the no-false-negatives condition of the oracle by construction.

A false positive occurs for an element $y \notin S$ if all $k$ bits corresponding to its hash values have been incidentally set to 1 by the insertion of other elements. Assuming the hash functions distribute elements uniformly over the bit array, the probability of a specific bit remaining 0 after the insertion of $n$ elements is $(1 - 1/m)^{kn} \approx e^{-kn/m}$. The probability of a false positive, $\varepsilon$, is the probability that all $k$ bits for the queried element are 1.

\begin{proposition}[False Positive Probability \cite{Tarkoma2012TheoryAP}]
    The false positive probability of a Bloom filter is $\varepsilon \approx (1 - e^{-kn/m})^k$.
\end{proposition}

A given ratio $m/n$ of bits per element determines an optimal number of hash functions that minimizes this probability.

\begin{theorem}[Optimal Hash Function Count \cite{Broder2004NetworkAO}]
    For a given ratio $m/n$, the false positive rate $\varepsilon$ is minimized when the number of hash functions is $k = (m/n) \ln 2$.
\end{theorem}

This theorem yields a direct relationship between the space allocated per element and the achievable false positive rate.

\begin{corollary}[Space-Error Relationship \cite{Broder2004NetworkAO}]
    To achieve a target false positive rate $\varepsilon$ with an optimal number of hash functions, the required number of bits per element is $m/n = -\frac{\ln \varepsilon}{(\ln 2)^2} \approx 1.44 \log_2(1/\varepsilon)$.
\end{corollary}

The asymptotic space complexity of the Bloom filter is therefore $\Theta(n \log(1/\varepsilon))$, independent of the size of the elements. The performance relies on high-quality hash functions. In practice, computationally inexpensive non-cryptographic hash functions are sufficient, and techniques such as double hashing, which generates $k$ hash values from two initial hashes, can reduce computational overhead without asymptotically increasing the false positive probability \cite{Tarkoma2012TheoryAP}.

The properties established above hold under two assumptions: the set $S$ is static after construction, and its cardinality $n$ is known when the filter is initialized. Distributed deployment introduces constraints that violate these assumptions. In systems where compute and storage are separated by network boundaries, where data undergoes continuous insertion and deletion across partitions, and where dataset cardinality is neither fixed nor predictable, the architectural requirements for membership oracles extend beyond the space-error trade-off. The subsequent analysis identifies these constraints and examines their implications for oracle design.

\subsection{The Constraint of State Synchronization}

In a distributed system where compute and storage are separated by network boundaries, a membership oracle must maintain coherence with a dataset that evolves through insertion and deletion operations across multiple partitions. This requirement for state synchronization introduces constraints not present in single-node deployments where the oracle and the data it represents are co-located. We identify two failure modes that emerge when a membership oracle cannot adapt its internal state to reflect changes in the underlying dataset.

The first failure mode is the persistence of stale positives. When an element is deleted from the dataset, a Bloom filter that represents that dataset cannot remove the element's contribution from its bit array. The bits set during the element's insertion remain set indefinitely. Subsequent queries for the deleted element will continue to return true. In a distributed key-value store that uses Bloom filters to avoid disk I/O for non-existent keys, this behavior forces the system to perform expensive read operations merely to discover that the key has been deleted or that a tombstone marker is present. The filter fails to converge to the true state of the system. The cost of this divergence is proportional to the delete rate of the workload; systems with high churn experience continuous wasted I/O as filters accumulate obsolete positive indications \cite{Mosharraf2021ImprovingLA, Tarkoma2012TheoryAP}.

The second failure mode is capacity exhaustion. A Bloom filter is parameterized for an expected cardinality $n$. When the actual number of inserted elements significantly exceeds this estimate, the bit array becomes saturated. As saturation increases, the false positive rate degrades toward unity, rendering the filter ineffective as a negative filter. In distributed environments where dataset cardinality is determined by aggregate insertion rates across multiple nodes and where traffic patterns are non-stationary, the capacity of a fixed-size filter becomes a constraint that requires either over-provisioning (wasting memory) or periodic reconstruction (introducing service interruption) \cite{Almeida2007ScalableBF, Tarkoma2012TheoryAP}.

\subsubsection{Architectural Variants and Their Trade-offs}

We analyze two architectural extensions to the standard Bloom filter that address the identified failure modes. Each extension resolves one constraint but introduces a new limitation that restricts its applicability in distributed contexts.

The Counting Bloom Filter addresses deletion by replacing the bit array with an array of counters. Each counter tracks the number of elements mapped to its position. Insertion increments the relevant counters; deletion decrements them. A query returns true if all corresponding counters are positive. This mechanism permits element removal without introducing false negatives for other elements.

\begin{definition}[Counting Bloom Filter \cite{Tarkoma2012TheoryAP}] \label{def:counting-bloom}
    A Counting Bloom Filter consists of an array $C$ of $m$ counters, each of $c$ bits, initialized to all zeros. It uses a collection of $k$ independent hash functions $\mathcal{H} = \{h_1, \ldots, h_k\}$, where each $h_i : \mathcal{U} \to \{0, \ldots, m-1\}$.
\end{definition}

The structure supports deletion but incurs a space overhead proportional to the counter width $c$. Analysis demonstrates that 4-bit counters suffice to prevent overflow for most workloads, yielding an approximately four-fold increase in memory consumption relative to a standard Bloom filter \cite{Fan1998SummaryCA, Tarkoma2012TheoryAP}. The arithmetic properties of counter arrays enable set difference operations through counter subtraction, a capability useful for distributed set reconciliation protocols \cite{Guo2013SetRV}.

The Scalable Bloom Filter addresses capacity exhaustion through a sequence of constituent filters with increasing sizes and decreasing error rates. When the active filter reaches a load threshold, a new, larger filter is appended to the sequence. Queries must inspect all filters in the sequence.

\begin{definition}[Scalable Bloom Filter \cite{Almeida2007ScalableBF}]
    An SBF is an ordered sequence of $s$ Bloom filters, $B_1, B_2, \ldots, B_s$. Each filter $B_i$ is defined by its size $m_i$ and number of hash functions $k_i$. The sequence is constructed such that for $i > 1$, $m_i > m_{i-1}$ and the target false positive rate $\varepsilon_i < \varepsilon_{i-1}$.
\end{definition}

This architecture accommodates unbounded growth but sacrifices a property critical for distributed aggregation. The heterogeneity of the constituent filters prevents bitwise operations such as the OR operation required for computing the union of two sets. Two SBF instances representing different sets cannot be merged through element-wise combination of their internal arrays because the arrays differ in size and hash function count. This absence of composability renders the SBF unsuitable for distributed query processing where partial results computed on different nodes must be aggregated \cite{Liu2014ParBFAP, Almeida2007ScalableBF}.

\subsubsection{The Requirement for Dynamic Mutability}

The analysis of the two failure modes and their proposed resolutions exposes a fundamental architectural requirement for membership oracles in distributed systems. An oracle that cannot delete elements accumulates stale state, forcing redundant operations that negate the performance benefits the oracle was designed to provide. An oracle that cannot scale capacity requires either over-provisioning or service interruption for reconstruction. The Counting Bloom Filter resolves deletion at the cost of substantial space overhead. The Scalable Bloom Filter resolves capacity at the cost of lost composability.

A membership oracle suitable for distributed deployment must satisfy three conditions: it must support element deletion in constant time, it must not require auxiliary data structures whose space overhead negates the oracle's compactness, and it must not introduce false negatives for elements that remain in the set. We formalize this requirement as the property of dynamic mutability.

\begin{definition}[Dynamic Mutability] \label{def:mutability}
    A probabilistic oracle $\mathcal{D}$ representing a set $S$ exhibits dynamic mutability if it supports both $\textsc{Insert}(x)$ and $\textsc{Delete}(x)$ operations for an element $x$ in $O(1)$ time complexity. The space complexity must increase by at most a constant multiplicative factor compared to a static probabilistic oracle with the same target false positive rate and expected capacity. The $\textsc{Delete}(x)$ operation must not introduce false negatives for any other element $y \in S, y \neq x$.
\end{definition}

This property is not a performance optimization but a structural requirement for oracles that must track the lifecycle of data in systems where deletion is a first-class operation. Distributed key-value stores that implement tombstones for deleted keys, Log-Structured Merge-Trees that compact data files and purge obsolete entries, and content distribution networks that expire cached objects all require membership oracles whose internal state can synchronize with these deletion events. The property of dynamic mutability characterizes the class of oracles that satisfy this requirement. We now examine the architectural realization of this property in the Cuckoo Filter.

\subsection{Cuckoo Filters: Architectural Realization of Dynamic Mutability} \label{subsec:cuckoo-filters}

The Cuckoo Filter satisfies the conditions of dynamic mutability established in Definition~\ref{def:mutability}. Its delete operation executes in $O(1)$ time, requires no auxiliary counters, and does not introduce false negatives for other elements. This property derives from its partial-key hashing scheme, which enables the relocation and removal of an element's fingerprint using only the fingerprint itself, decoupling the operation from access to the original element. We analyze the technical mechanisms that provide this capability and examine the distributed deployment patterns that dynamic mutability enables.

\subsubsection{Partial-Key Hashing} \label{subsubsec:partial-key-hashing}

The fundamental challenge in deploying Cuckoo Filters across network boundaries lies in the relocation operation itself. During insertion, when both candidate buckets are full, the filter must displace an existing fingerprint to its alternate location. Standard cuckoo hashing accomplishes this by recomputing the hash of the original element; the element's full key must be available to determine where to move it. In a distributed system, however, filters may be partitioned or replicated across nodes, and the original element may not be co-located with the filter metadata. Retrieving the full element for relocation decisions across network boundaries would introduce prohibitive overhead.

The solution is partial-key cuckoo hashing, introduced by Fan et al., which enables relocation decisions using only the stored fingerprint \cite{Fan2014CuckooFP}. This capability is the prerequisite for all subsequent distributed patterns.

\begin{definition}[Partial-Key Cuckoo Hashing \cite{Fan2014CuckooFP}]
    For an element $x \in \mathcal{U}$ with precomputed fingerprint $f = \textsc{fingerprint}(x)$, the two candidate bucket indices are computed as:
    \begin{equation*}
        i_1 = \textsc{hash}(x), \quad i_2 = i_1 \oplus \textsc{hash}(f)
    \end{equation*}
    where $\oplus$ denotes bitwise XOR. There is a symmetry that allows us to, given $i_2$ and $f$, recover $i_1$ via the same equation: $i_1 = i_2 \oplus \textsc{hash}(f)$.
\end{definition}

This symmetry decouples relocation from full-key access. When a fingerprint must be moved from bucket $i_1$ to its alternate location, the system computes the destination using only $f$ and the current position. The fingerprint can traverse network boundaries, be temporarily replicated on different nodes, or be moved as part of partition rebalancing, all without requiring access to the original element. For distributed systems where filters may be stored separately from data, or where data undergoes continuous replication and migration, this property is essential.

\begin{definition}[Cuckoo Filter Core Operations \cite{Fan2014CuckooFP}]
    A Cuckoo Filter with $m$ buckets, each containing $b$ entries, supports three operations:

    \emph{Lookup}$(x)$: Compute $i_1 = \textsc{hash}(x)$, $i_2 = i_1 \oplus \textsc{hash}(f)$. Return true if fingerprint $f$ exists in either bucket $i_1$ or $i_2$; otherwise return false. Query always examines at most $2b$ entries.

    \emph{Insert}$(x)$: Compute $i_1, i_2$. If either bucket contains an empty entry, place $f$ there and complete. Otherwise, select one bucket randomly, displace its resident fingerprint $f'$, compute the alternate location for $f'$ using $i \gets i \oplus \textsc{hash}(f')$, and repeat the displacement process. If relocations exceed a configurable threshold (typically 500), insertion fails and a capacity increase is required.

    \emph{Delete}$(x)$: Compute $i_1, i_2$. If $f$ exists in either bucket, remove one copy; otherwise return failure. Deletion requires no counters or additional bookkeeping.
\end{definition}

We analyze the probability of a false positive, which occurs when an element not in the set hashes to a fingerprint that collides with an existing fingerprint in one of its two candidate buckets. A query inspects at most $2b$ entries across two buckets. This inspection process yields a direct upper bound on the false positive probability, which we state formally.

\begin{proposition}[False Positive Probability Bound \cite{Fan2014CuckooFP}]
    Let a Cuckoo Filter be configured with $b$ entries per bucket and a fingerprint length of $f$ bits. The false positive probability $\varepsilon$ satisfies the inequality:
    \begin{equation*}
        \varepsilon \leq \frac{2b}{2^f}
    \end{equation*}
\end{proposition}

From this inequality, we derive the minimum fingerprint size $f$ required to achieve a target false positive probability $\varepsilon$:
\begin{equation*}
    f \geq \lceil \log_2(1/\varepsilon) + \log_2(2b) \rceil
\end{equation*}
For a configuration with $b=4$ and a target $\varepsilon=0.01$ (a 1\% false positive rate), this relationship requires a fingerprint size of at least $f \geq \lceil \log_2(100) + \log_2(8) \rceil = \lceil 6.64 + 3 \rceil = 10$ bits. This yields a compact representation that is computationally sufficient to mitigate a high rate of collisions.

The architectural consequences of this design become evident when we consider the delete operation. A standard Bloom filter offers no delete mechanism; removing an element's bits risks corrupting the representation of other stored elements. A Counting Bloom Filter enables deletion but requires per-position counters, which results in a space overhead of approximately four times that of a standard Bloom filter \cite{Fan2014CuckooFP}. In contrast, the Cuckoo Filter implements deletion by locating the specified fingerprint in one of its two candidate buckets and removing it. This operation requires no auxiliary data structures and directly enables the dynamic distributed architectures we analyze subsequently.

\subsubsection{Capacity Scaling}

In a distributed system, the cardinality of a dataset is not known \emph{a priori}, requiring any data structure to accommodate significant growth. A Cuckoo Filter of a fixed size cannot accept new insertions once its load factor approaches its theoretical limit. While single-node deployments address this by rebuilding the filter into a larger table offline, such a procedure is untenable in distributed systems that operate under strict availability and latency constraints. A node that halts operations to rebuild its local filter violates service agreements and introduces a point of coordination that undermines distributed design principles.

The Dynamic Cuckoo Filter (DCF), introduced by Chen et al., presents a systematic solution to this scaling problem \cite{Chen2017TheDC}. Instead of rebuilding a full filter, the DCF architecture allocates a new, empty filter and appends it to an ordered sequence of filters. All new insertions are directed to the most recently appended filter.

\begin{definition}[Dynamic Cuckoo Filter \cite{Chen2017TheDC}]
    A DCF is an ordered sequence of Cuckoo Filters, $\text{CF}_1, \text{CF}_2, \ldots, \text{CF}_k$. A pointer, $\text{curCF}$, references the terminal filter in the sequence. An insertion operation first attempts placement on $\text{curCF}$. If this insertion fails due to capacity limits, a new filter is allocated, appended to the sequence, and designated as the new $\text{curCF}$. A query operation must traverse this sequence in reverse chronological order, from $\text{curCF}$ to $\text{CF}_1$, until the target fingerprint is found or all filters have been examined.
\end{definition}

The DCF architecture permits unbounded capacity growth without service interruption, but it achieves this at the cost of query performance. A query for an element $x$ requires a sequential search through the filter chain. If $x$ was inserted into an early filter and the chain has subsequently grown to a length of $k$, the query must inspect all $k$ filters.

\begin{proposition}[DCF Query Complexity]
    A membership query in a DCF composed of $k$ constituent filters requires $O(k)$ bucket accesses in the worst case. For a system containing $n$ total items, where each filter maintains a load factor $\alpha \approx 0.95$ and has a capacity of $C$ items, the number of filters scales as $k = \Theta(n/C)$. Consequently, query latency grows proportionally with the total dataset size.
\end{proposition}

This architectural pattern creates a direct conflict between capacity scaling and query performance. The mechanism that enables system growth is the same mechanism that causes query latency to degrade. For modern distributed systems where query latency is a critical performance metric, such as in data serving, caching, and real-time analytics, a linear growth in query time is architecturally unacceptable.

The linear query complexity arises because the DCF organizes its constituent filters temporally, by insertion order, rather than deterministically by element hash value. No direct mapping exists between an element's content and the specific filter that contains it. Membership testing therefore degenerates to a sequential search, an operation incompatible with the constant-time query requirements of scalable distributed systems.

A scalable architecture requires that the mapping from an element to its storage location be a deterministic function of the element's hash value, independent of the system's growth over time or the insertion order of other elements. This decouples the problem of element location from the problem of capacity management, ensuring that query performance remains constant as the system scales.

\subsubsection{Consistent Hashing}

The linear degradation in query performance of the Dynamic Cuckoo Filter arises from its temporal organization. We observe that distributed systems have already addressed analogous scaling challenges for data partitioning. Architectures such as Apache Cassandra and Amazon DynamoDB employ consistent hashing to distribute their keyspace \cite{Lakshman2010CassandraAD, DeCandia2007DynamoAH}. In these systems, nodes are assigned token ranges on a virtual hash ring. An element's hash value deterministically maps it to a specific point on the ring, and thus to a responsible node. The addition or removal of a node requires remapping only a localized subset of the keyspace, preserving a stable and deterministic mapping for the majority of data. Redis Cluster achieves a similar outcome using a hash slot partitioning scheme, where a fixed number of slots are distributed among nodes and can be migrated individually to scale the cluster \cite{redis-cluster-spec}.

This same principle can be applied to the organization of a Cuckoo Filter's buckets. Instead of organizing filters in a temporal sequence, we can distribute the buckets themselves across a consistent hash ring. This architectural shift aligns the filter's structure with established distributed partitioning schemes.

\begin{definition}[Index-Independent Cuckoo Filter \cite{Luo2021ACC}]
    An Index-Independent Cuckoo Filter (I2CF) maintains its buckets on a consistent hash ring. For an element $x$ with fingerprint $f$, its candidate bucket locations are determined by hashing the element to points on this ring. The system identifies the first bucket successor to each of these points on the ring as a candidate bucket.
\end{definition}

The defining property of the I2CF is that bucket identifiers depend only on the element's hash value, not on the total number of buckets $m$ in the system. Capacity expansion is achieved by adding a new bucket to a chosen position on the ring. This operation necessitates relocating only those elements whose hash values fall within the range immediately preceding the new bucket's position. For a system with $n$ items distributed across $m$ buckets, a rebalancing event affects an expected $O(n/m)$ items, a small fraction of the total dataset.

This design decouples capacity growth from query complexity, resolving the architectural conflict present in the DCF.

\begin{theorem}[I2CF Query Complexity]
    For an I2CF with $m$ buckets organized on a consistent hash ring, a membership query for an element $x$ requires a sequence of constant-time operations: (1) hash computations to determine ring positions, (2) a ring lookup to identify candidate buckets (achievable in $O(1)$ expected time with an appropriate auxiliary index), and (3) inspection of entries within those buckets. The total query complexity is therefore $O(1)$, independent of the total number of items $n$ and the total number of buckets $m$.
\end{theorem}

The I2CF architecture resolves the conflation of capacity management and element mapping found in the DCF. The DCF modifies the element-to-filter mapping to accommodate growth, leading to performance degradation. In contrast, the I2CF maintains a fixed, deterministic mapping from elements to ring positions while allowing the set of buckets to expand elastically. As a result, growth and query performance become orthogonal concerns.

\subsubsection{Co-Partitioning Filters with Distributed Data} \label{subsubsec:co-partitioning}

We observe that the architectural value of the Index-Independent Cuckoo Filter is most fully realized when its partitioning scheme is aligned with that of the host distributed system. Many distributed key-value stores partition their keyspace across a set of nodes using a consistent hash ring. In such systems, each node assumes responsibility for one or more token ranges on the ring.

We apply this same partitioning logic to the filter's buckets, an approach whose efficacy is demonstrated by Luo et al. in their work on distributed filters over Redis Cluster \cite{Luo2018DistributedDC}. We formalize this alignment as a co-partitioned filter architecture.

\begin{definition}[Co-Partitioned Filter Architecture \cite{Luo2018DistributedDC}]
    A distributed system of $m$ nodes maintains a consistent hash ring, where each node $n_i$ is assigned responsibility for a token range $[t_i, t_{i+1})$. The node maintains: (1) all data items $d$ such that $H(d) \in [t_i, t_{i+1})$, and (2) all Cuckoo Filter buckets $b$ such that $\text{position}(b) \in [t_i, t_{i+1})$. A membership query for an element $x$ is routed to the node responsible for the token range containing $H(x)$ for local processing.
\end{definition}

This architecture instantiates the co-partitioning pattern identified in the framework (Section~\ref{subsubsec:mutability-copartition}). The pattern emerges as a consequence of dynamic mutability: the ability to execute deletion locally permits the stable assignment of filter buckets to consistent hash ranges, as element removal does not require global state coordination. A query processed on the node that owns the element's token range incurs only the cost of local memory access. A query originating from a different node requires a single network round-trip. The latency difference between local memory access (microseconds) and network communication (milliseconds) is typically three to four orders of magnitude.

We can model the expected query latency for this architecture under a uniform query distribution. In a system with $m$ nodes, a randomly routed query has a probability $1/m$ of being processed locally and a probability $(1 - 1/m)$ of requiring network transit. This leads to the following formulation for the expected query latency.

\begin{theorem}[Co-Partitioned Query Latency] \label{thm:co-partitioned-latency}
    In a co-partitioned I2CF architecture of $m$ nodes under a uniform query distribution, the expected query latency $\mathbb{E}[L_q]$ is given by:
    \begin{equation*}
        \mathbb{E}[L_q] = \frac{1}{m} t_{\text{local}} + \left(1 - \frac{1}{m}\right) t_{\text{network}}
    \end{equation*}
    where $t_{\text{local}}$ represents local bucket access time and $t_{\text{network}}$ represents network round-trip latency.
\end{theorem}

This latency bound is independent of the total dataset size $n$. Query performance is determined by hardware characteristics and system scale ($m$), not by the number of elements stored.

The alignment creates a synergy with the underlying infrastructure of many distributed systems. By reusing an existing consistent hashing ring, the system gains a scalable membership testing capability with minimal additional coordination overhead. Fault tolerance mechanisms, such as data replication, can be extended to include filter buckets, allowing replicated filter state to propagate through existing replication channels.

The experimental results in \cite{Luo2018DistributedDC} confirm the efficacy of this approach, demonstrating significant memory savings over Bloom filter alternatives, full support for deletions, and seamless capacity scaling as nodes are added to the cluster.

\subsubsection{Hierarchical Elasticity}

While the I2CF architecture provides fine-grained, bucket-level elasticity, its rebalancing mechanism is most effective under predictable, incremental growth in dataset cardinality. Distributed systems, however, often experience non-uniform growth patterns, including sudden spikes in insertion rates. In such scenarios, where the rate of cardinality growth exceeds the rate at which new buckets can be added and populated, the system may experience cascading insertion failures.

To address this class of workloads, we analyze the Consistent Cuckoo Filter (CCF), an architecture that introduces a second, coarser tier of elasticity by organizing a collection of independent I2CF instances \cite{Luo2021ACC}.

\begin{definition}[Consistent Cuckoo Filter \cite{Luo2021ACC}]
    A CCF is a collection of I2CF instances, $\{I_1, I_2, \ldots, I_s\}$, initially with $s=1$. Insertions are directed to a designated active instance, typically $I_s$. When $I_s$ reaches a predefined load factor threshold, a new, empty instance $I_{s+1}$ is allocated and becomes the active target for subsequent insertions. A query operation must inspect all instances in the collection until the target fingerprint is found.
\end{definition}

This architecture allows the system to increase its total capacity instantly by adding a new, untapped I2CF. However, this "scale-out" mechanism, if left unmanaged, structurally resembles the chaining model of the Dynamic Cuckoo Filter. A query must probe every instance in the collection, leading to a performance dependency on the number of instances, $s$.

\begin{proposition}[CCF Unmanaged Query Complexity]
    In a CCF where new, fixed-capacity I2CF instances are added sequentially without any consolidation, the number of instances $s$ grows linearly with the total number of items $n$. The worst-case query complexity is therefore $O(s)$, which is equivalent to $O(n)$.
\end{proposition}

\begin{proof}
    Let $n$ be the total number of items and let each I2CF instance have an effective capacity of $C$ items (e.g., $C = m \cdot b \cdot \alpha$, where $m$ is the number of buckets, $b$ is the entries per bucket, and $\alpha$ is the target load factor). In an unmanaged growth model, a new instance is added every time the system has stored an additional $C$ items. The number of instances $s$ required to store $n$ items is therefore $s = \lceil n/C \rceil$. As $n \to \infty$, this gives $s = \Theta(n)$. Since a query must inspect all $s$ instances in the worst case, the query complexity is $O(s) = O(n)$.
\end{proof}

The CCF architecture as defined by Luo et al. explicitly includes management policies to prevent this linear degradation \cite{Luo2021ACC}. It is not merely a chain of filters but a managed system that provides hierarchical scaling. This system operates at two granularities. The first is a fine-grained, bucket-level elasticity, where capacity adjustments occur within each constituent I2CF by adding or removing buckets from its consistent hash ring. The second is a coarse-grained, filter-level elasticity, which accommodates sudden load spikes by adding entire new I2CF instances. To control the growth of query latency, the architecture incorporates a consolidation mechanism, which periodically merges under-utilized instances to reduce the total number of instances, $s$. This hierarchical approach allows the system to manage capacity at multiple granularities, handling gradual growth through the fine-grained mechanism and absorbing sudden traffic spikes by adding new instances, while maintaining long-term performance through background consolidation.

\subsubsection{Concurrency Control Strategies}

The deployment of Cuckoo Filters in a distributed environment introduces concurrency hazards. During an insertion operation, the filter performs a sequence of displacements that modify multiple bucket locations. A concurrent query examining these locations during this transient state may fail to find an element that has been evicted from its original bucket but not yet placed in its alternate location. This race condition, which we term a \emph{false miss}, results in a definitive negative answer for an element present in the set, thereby violating the data structure's membership guarantee.

We analyze three distinct concurrency control strategies, each designed to resolve the false-miss problem for a specific system architecture and workload pattern.

The first strategy, optimistic concurrency, is designed for read-dominant workloads, such as those found in caching systems. Fan et al. developed this model for the MemC3 system \cite{Fan2013MemC3CA}. The model serializes write operations while allowing multiple readers to proceed without locks, using version counters for synchronization. A key component of this model is the execution order of relocations along a displacement chain. After identifying a valid path, the system moves the last element in the chain into the empty slot first, which frees a new slot. The second-to-last element is then moved into this newly vacated slot. This reverse-order execution ensures that every element being relocated remains accessible in at least one of its two candidate buckets at all times.

\begin{definition}[Optimistic Cuckoo Hashing \cite{Fan2013MemC3CA}]
    A writer maintains a version counter for each bucket. Before initiating relocations, it identifies a valid displacement path (e.g., $a \to b \to c \to \text{empty}$) without acquiring locks. The relocations are then executed in reverse order, with each move incrementing the version counter of the involved bucket. A reader records the versions of its two candidate buckets before inspection and re-checks them afterward. If any version has changed, the reader discards its result and retries the operation.
\end{definition}

This strategy is highly effective for workloads where the read-to-write ratio is high, as readers do not incur lock acquisition overhead.

For general-purpose multi-core systems with more balanced read-write workloads, we analyze a fine-grained locking model that enables greater writer parallelism. The libcuckoo library\footnote{\url{https://github.com/efficient/libcuckoo}} implements this model using lock striping, where an array of independent locks protects disjoint sets of buckets \cite{Li2014AlgorithmicIF, Li2023LockFreeBC}. A thread acquires locks only for the buckets involved in its operation, which permits concurrent write operations on non-overlapping regions of the filter.

For disaggregated memory architectures that utilize Remote Direct Memory Access (RDMA), in-memory locking mechanisms are inapplicable. Clients access remote memory directly over the network. Grant et al. developed RCuckoo for this environment, which employs a lease-based locking protocol \cite{Grant2025CuckooFC}. In this model, clients acquire time-bounded leases on remote buckets via atomic RDMA operations. The automatic expiration of these leases provides fault tolerance to client failures without requiring explicit recovery protocols.

\begin{definition}[Lease-Based Locking for RDMA \cite{Grant2025CuckooFC}]
    For disaggregated memory architectures, we manage client failures during write operations through a lease-based protocol. The system maintains a separate lease table in remote memory, which is distinct from the primary lock table for normal operations. A client identifies a potentially stranded lock when an operation exceeds a predefined timeout threshold. To initiate recovery, the client acquires an exclusive repair lease for the affected index region via an atomic compare-and-swap operation on the lease table. This operation grants the client exclusive rights to modify the table state for recovery purposes. Upon completion of the repair, the client relinquishes the lease by clearing its ownership flag. The repair lease itself is subject to a timeout mechanism, which provides fault tolerance against clients that fail during the recovery process itself.
\end{definition}

We note that across all three concurrency models, the no-false-negatives property of the Cuckoo Filter remains invariant. A query that returns false provides a correct answer: the element is not, and has never been (unless subsequently deleted), a member of the set. The specific concurrency strategy is an implementation detail dictated by the hardware topology and workload, but the membership guarantee is preserved.

\subsubsection{Replication and Convergence}

When we replicate Cuckoo Filters across multiple nodes for fault tolerance, the replicas may temporarily diverge due to network latency in propagating updates. We analyze the consequences of this divergence, particularly for deletion operations, and compare the behavior of Cuckoo Filters to that of standard and Counting Bloom Filters.

A standard Bloom filter provides no mechanism for deletion; elements, once inserted, remain in the filter indefinitely. A Counting Bloom Filter supports deletion by decrementing counters but at the cost of an approximately four-fold increase in space overhead \cite{Fan2014CuckooFP}. A Cuckoo Filter, by contrast, implements deletion as a direct removal of the element's fingerprint.

In a replicated environment, this distinction becomes architecturally significant. Consider a Cuckoo Filter replicated on nodes $n_1$ and $n_2$. An insertion of element $x$ on $n_1$ propagates asynchronously to $n_2$. During the propagation window, a query for $x$ on $n_2$ will return false. This behavior, a form of temporal false negative, is inherent to any asynchronously replicated system.

The key difference emerges in deletion scenarios, a problem central to systems that use tombstones to mark deleted data under eventual consistency \cite{Mosharraf2021ImprovingLA}. When element $x$ is deleted on $n_1$, its fingerprint is immediately removed from the local filter. During the propagation window, a query for $x$ to $n_1$ correctly returns false, while a query to $n_2$ returns a stale positive. Once the deletion operation propagates to $n_2$, its fingerprint is also removed, and both nodes converge to a consistent state where $x$ is correctly identified as absent.

A Bloom filter behaves differently. Since it cannot remove the element, it will continue to return positive for the deleted key indefinitely. This forces the system to perform expensive validation lookups against the backing data store merely to discover a tombstone, negating much of the filter's performance benefit in workloads with frequent deletions \cite{Mosharraf2021ImprovingLA}.

\begin{proposition}[Cuckoo Filter Replication Convergence]
    In an eventually consistent system with replicated Cuckoo Filters, let an element $x$ be inserted at time $t_i$ and deleted at time $t_d > t_i$. For any replica $r$, there exist propagation delays such that the insertion is applied at time $\tau_i \geq t_i$ and the deletion is applied at time $\tau_d \geq t_d$. For all times $t > \tau_d$, a query for $x$ on replica $r$ will deterministically return false. In contrast, a standard Bloom filter provides no upper bound on the time until a query for a deleted element ceases to return a positive result.
\end{proposition}

This property enables Cuckoo Filters to converge more rapidly to the true state of the system, especially in environments with frequent deletions. For distributed systems implementing eventual consistency, this allows the filter's state to accurately track operations such as the purging of tombstones during compaction or garbage collection. The ability to remove fingerprints, rather than merely marking them for future cleanup, ensures that the filter remains a reliable oracle for the live, non-deleted state of the dataset.

\subsubsection{Impact on Distributed System Architectures} \label{subsubsec:arch-impact}

We analyze the concrete architectural impact of Cuckoo Filters in several domains where distributed systems require both efficient membership testing and dynamic element deletion.

Our first application domain is Log-Structured Merge-Tree (LSM-tree) based key-value stores. In these systems, data is organized into immutable, sorted files on disk (SSTables), each typically indexed by an in-memory filter. Deletions are often handled by writing a tombstone marker rather than by immediate data removal \cite{Mosharraf2021ImprovingLA}. A standard Bloom filter, being unable to remove entries, continues to return a positive match for a deleted key. This behavior forces the database to perform a disk I/O operation only to discover the tombstone, incurring unnecessary latency. A Cuckoo Filter resolves this inefficiency. During data compaction, the fingerprint corresponding to a purged tombstone can be removed from the filter. Subsequent queries for the deleted key will correctly return false, eliminating the wasted I/O. The Chucky architecture further exploits this by unifying the multiple per-level filters of an LSM-tree into a single Cuckoo Filter, which reduces the point query complexity from $O(\text{number of levels})$ to $O(1)$ \cite{Dayan2021ChuckyAS}.

In named-data networking (NDN), we examine the Pending Interest Table (PIT), which tracks outstanding content requests. The DiCuPIT system utilizes distributed Cuckoo Filters to replace traditional Bloom filter-based implementations. This architectural change yields a 36\% reduction in lookup time and a 31\% reduction in memory consumption compared to Bloom filter baselines \cite{Mahmoudi2023DiCuPITDC}. The primary advantage is the ability to correctly and efficiently delete PIT entries as their corresponding data responses arrive, a fundamental operation that standard Bloom filters cannot support.

Our third domain is state synchronization in blockchain networks. Nodes in these systems maintain a mempool of pending transactions. To propagate a new block efficiently, the Gauze protocol employs Cuckoo Filters to summarize the contents of a node's mempool \cite{ding2023gauze}. A node propagating a block sends this compact filter to its peers. The receiving nodes use the filter to identify which transactions in the new block are absent from their local mempool and request only this delta. This method reduces block propagation bandwidth by 40 to 50\%. The efficiency gain is a direct result of the filter's compact representation; a filter for a million-element set requires only a few megabytes of space.

These applications demonstrate a shared architectural requirement: efficient support for both insertion and deletion in resource-constrained environments, where the primary bottlenecks are network bandwidth, I/O latency, or processing throughput. The ability to delete elements distinguishes the Cuckoo Filter from the Bloom filter not as an incremental improvement, but as an enabling architectural property. It facilitates correct state reconciliation, efficient rebalancing under dynamic scaling, and accurate membership semantics in systems where the lifecycle of data includes explicit deletion.

\subsection{Distributed Patterns and the Mutability-Composability Tradeoff}

The analysis in subsection~\ref{subsec:bloom-limitations} established that the standard Bloom filter lacks dynamic mutability: it cannot delete elements without introducing false negatives or incurring substantial space overhead through auxiliary counters. The variants presented (Counting Bloom Filter, Scalable Bloom Filter) each resolve one constraint while introducing another. The Counting Bloom Filter provides deletion at a four-fold space cost. The Scalable Bloom Filter provides capacity growth but sacrifices the algebraic properties required for distributed aggregation.

We now examine distributed deployment patterns developed specifically for Bloom filters. These patterns address challenges that emerge when filters are replicated or partitioned across nodes: temporal inconsistency, heterogeneous filter parameters, and the need for state aggregation. We analyze two representative architectures and classify them according to the framework properties they provide. This classification exposes a structural limitation: the patterns address specific deployment challenges but cannot overcome the absence of dynamic mutability in the underlying data structure.

\subsubsection{Parallel Partitioned Bloom Filter}

The Scalable Bloom Filter, documented in subsection~\ref{subsec:bloom-limitations}, accommodates unbounded growth through a sequence of filters with heterogeneous parameters. This heterogeneity prevents bitwise operations such as OR for set union. Two SBF instances cannot be merged because their constituent filters differ in size and hash function count. This absence of composability renders the SBF unsuitable for distributed aggregation tasks where partial results from different nodes must be combined.

The Parallel Partitioned Bloom Filter (ParBF) addresses this limitation by constructing a homogeneous sequence of sub-filters \cite{Liu2014ParBFAP}. Each sub-filter in the sequence shares the same size $m$ and the same number of hash functions $k$. This structural uniformity preserves the algebraic properties required for distributed operations.

\begin{definition}[Parallel Partitioned Bloom Filter \cite{Liu2014ParBFAP}]
    A ParBF is a sequence of $s$ homogeneous Bloom filters, $B_1, B_2, \ldots, B_s$. Each filter $B_i$ has size $m$ and uses the same $k$ hash functions. New elements are inserted into the filter with the lowest current load. When all filters exceed a threshold load factor, a new filter $B_{s+1}$ with the same parameters $(m, k)$ is appended to the sequence.
\end{definition}

The homogeneity of the constituent filters enables bitwise union operations. Two ParBF instances representing sets $S_1$ and $S_2$ can be merged by computing the bitwise OR of corresponding sub-filters: the $i$-th sub-filter of the merged ParBF is the OR of the $i$-th sub-filters from both inputs. This operation produces a ParBF representing the set union $S_1 \cup S_2$. The result satisfies the one-sided error property: queries for elements in $S_1 \cup S_2$ return true with probability 1, and queries for elements not in the union return false positive with a probability bounded by the false positive rates of the constituent filters.

The architecture further partitions the sequence of sub-filters into groups to enable parallel query execution. A query operation can be distributed across multiple threads, each inspecting a disjoint subset of the filter sequence. This parallelism improves throughput in high-concurrency systems.

We evaluate the ParBF through the lens of the framework established in Section~\ref{sec:framework}. The structure provides a form of algebraic composability: the bitwise OR operation is commutative, associative, and idempotent. Two ParBF instances can be merged without coordination regarding the order of operations. However, the ParBF does not satisfy the conditions of dynamic mutability as defined in Definition~\ref{def:mutability}. The underlying Bloom filters cannot delete elements. Once an element is inserted into any sub-filter in the sequence, its contribution to the bit array is permanent. A delete operation would require either clearing bits (introducing false negatives) or maintaining auxiliary counters (violating the space complexity constraint).

The ParBF represents an architectural solution to the composability limitations of the Scalable Bloom Filter. It enables distributed aggregation through bitwise operations but inherits the fundamental constraint of the standard Bloom filter: the inability to synchronize its state with a dataset that undergoes deletion.

\subsubsection{Distributed Bloom Filter Protocol}

In a replicated system where Bloom filters are maintained on multiple nodes, asynchronous propagation of updates introduces temporal inconsistency. As documented in subsection~\ref{subsubsec:temporal-fn}, when an element is inserted on node $n_1$ and the update has not yet propagated to replica $n_2$, a query for that element on $n_2$ will return false. This behavior constitutes a temporal false negative: from the perspective of the global system state, the element is present, yet the oracle returns a definitive negative.

The Distributed Bloom Filter (DBF) protocol, introduced by Ramabaja et al., addresses this inconsistency through a probabilistic convergence mechanism \cite{Ramabaja2019TheDB}. The protocol does not attempt to eliminate temporal false negatives during the propagation window. Instead, it structures peer-to-peer interactions to guarantee convergence over multiple synchronization rounds while minimizing bandwidth consumption.

The protocol employs unique, pair-wise hash functions for each interaction between nodes. This design ensures that hash collisions are statistically independent across successive synchronization rounds. Each node pair derives a unique hash function family from a seed computed as the XOR of their node identifiers.

\begin{definition}[Distributed Bloom Filter Protocol \cite{Ramabaja2019TheDB}]
    Let two nodes, $n_i$ and $n_j$, with unique identifiers $ID_i$ and $ID_j$, wish to reconcile their local sets. The DBF protocol employs a unique, pair-wise hash function family $\mathcal{H}_{ij}$ generated from a seed $s_{ij}$, typically computed as $s_{ij} = ID_i \oplus ID_j$. A set of $k$ intermediate hashes $\{h_1^{ij}, \ldots, h_k^{ij}\}$ is derived from this seed. For each element $x \in \mathcal{U}$ with a pre-computed hash $H(x)$, the final hash values are computed as:
    \begin{equation*}
        H_l^{ij}(x) = (H(x) \oplus h_l^{ij}) \pmod m \quad \forall l \in \{1, \ldots, k\}
    \end{equation*}
    This method allows for the computationally inexpensive generation of unique Bloom filters for each peer interaction.
\end{definition}

The protocol deliberately utilizes filters with a high individual false positive rate (e.g., 50\%) to minimize bandwidth consumption. While a single exchange is highly probabilistic, the use of unique mappings for each interaction ensures that true set differences are revealed consistently, while random errors from false positives average out over multiple interactions. This approach sacrifices single-shot accuracy for rapid long-term convergence, a trade-off well-suited for eventually consistent systems \cite{Ramabaja2019TheDB}. The system-wide error rate diminishes exponentially with the number of participating nodes.

\begin{proposition}[DBF System-Wide Error Rate \cite{Ramabaja2019TheDB}]
    Let a system of $N$ nodes use the DBF protocol, where each individual filter has a false positive rate $\varepsilon$. For an element $y \notin \bigcup_{i=1}^N S_i$, the probability that a query for $y$ from a specific node to all of its $N-1$ peers results in a system-wide false positive is $\varepsilon_{\text{System}} \leq \varepsilon^{N-1}$.
\end{proposition}
\begin{proof}
    Let $\varepsilon = (1 - e^{-kn/m})^k$ be the FPR for a single filter. A system-wide false positive for a query from node $n_j$ regarding an element $y$ occurs if all $N-1$ peers it communicates with independently return a false positive. Let $E_i$ be the event of a false positive from peer $n_i$. The DBF protocol's use of unique seeds $s_{ji}$ for each peer interaction ensures that the hash function families $\mathcal{H}_{ji}$ are statistically independent for each peer $i$. Consequently, the events $E_i$ are independent. The probability of a joint failure is the product of their individual probabilities: $P(\bigcap_{i=1}^{N-1} E_i) = \prod_{i=1}^{N-1} P(E_i) = \varepsilon^{N-1}$.
\end{proof}

We evaluate the DBF protocol through the framework lens. The protocol addresses temporal inconsistency through probabilistic convergence but does not provide dynamic mutability. The underlying Bloom filters cannot delete elements; once an element is inserted during a synchronization round, it remains in the local filter indefinitely. If an element is deleted from the true dataset on some node, the protocol has no mechanism to propagate this deletion to peer filters.

The protocol also does not provide algebraic composability in the sense defined in Definition~\ref{def:composability}. While the protocol uses bitwise OR operations to combine filters during synchronization rounds, the use of unique hash functions for each node pair means that filters constructed on different nodes are not directly mergeable. Two filters representing the same set but constructed with different hash seeds cannot be combined through a bitwise OR to produce a valid filter for the union. The protocol achieves eventual consistency through iterative pairwise reconciliation, not through a single, order-independent merge operation.

The DBF protocol represents an architectural solution to the temporal inconsistency problem in replicated Bloom filters. It provides probabilistic guarantees about convergence rates and system-wide error bounds. However, it operates within the constraints of the underlying data structure: filters that cannot delete elements and that require coordination-specific hash functions for each interaction.

\subsubsection{Classification and Architectural Implications}

We classify the distributed patterns according to the framework properties established in Section~\ref{sec:framework}. The Parallel Partitioned Bloom Filter provides a restricted form of algebraic composability. The homogeneity of its constituent filters enables bitwise union operations that are commutative, associative, and idempotent. This property allows ParBF instances to be merged without coordination, supporting distributed aggregation tasks. However, the ParBF does not provide dynamic mutability. It inherits the deletion constraint of the standard Bloom filter.

The Distributed Bloom Filter protocol addresses temporal inconsistency through probabilistic convergence. It provides guarantees about the rate at which replicas converge to a consistent view of the dataset and bounds the system-wide false positive rate. However, the protocol does not provide dynamic mutability: filters cannot delete elements. It also does not provide algebraic composability in the framework sense: filters constructed with different hash seeds cannot be merged through a single, deterministic operation.

The distributed patterns developed for Bloom filters address specific deployment challenges but do not resolve the fundamental limitation: the absence of dynamic mutability. These patterns represent architectural workarounds rather than structural solutions. The ParBF works around the heterogeneity problem of the Scalable Bloom Filter to enable aggregation, but it cannot synchronize filter state with a dataset that undergoes deletion. The DBF protocol works around temporal inconsistency through iterative reconciliation, but it cannot remove obsolete entries from filters.

The Cuckoo Filter, by satisfying the conditions of dynamic mutability as established in subsection~\ref{subsec:cuckoo-filters}, enables a different class of distributed architectures. The co-partitioning pattern documented in subsection~\ref{subsubsec:co-partitioning} relies on the ability to execute deletion locally without global coordination. The replication convergence analysis in subsection~\ref{subsec:cuckoo-filters} demonstrates that Cuckoo Filter replicas converge to the true state of the dataset after deletions propagate, a property that Bloom filter replicas cannot achieve. The architectural impact case studies in subsection~\ref{subsubsec:arch-impact} document domains where the ability to delete elements is not an optimization but a requirement: LSM-tree compaction, named-data networking PIT management, and blockchain mempool synchronization.

The framework dichotomy between dynamic mutability and algebraic composability, introduced in Section~\ref{sec:framework}, explains why Bloom filter patterns and Cuckoo Filter patterns address different architectural contexts. Bloom filter patterns optimize for aggregation (ParBF) or eventual consistency (DBF) in scenarios where the dataset is append-only or where deletion is rare enough that stale positives are acceptable. Cuckoo Filter patterns optimize for state synchronization in scenarios where data undergoes explicit deletion and where the filter must accurately reflect the live, non-deleted state of the dataset. The choice between these pattern classes is dictated by the operational requirements of the distributed system, not by a comparison of space-error trade-offs.
